\documentclass[10pt,a4paper]{article}
\usepackage[margin=15mm,top=14mm,bottom=14mm]{geometry}
\usepackage[T1]{fontenc}
\usepackage[utf8]{inputenc}
\usepackage{lmodern}
\usepackage[table]{xcolor}
\usepackage{tabularx}
\usepackage{enumitem}
\usepackage{microtype}
\usepackage{hyperref}
\usepackage{titlesec}
\usepackage{fancyhdr}

\definecolor{ink}{HTML}{17352E}
\definecolor{green}{HTML}{146B4F}
\definecolor{mint}{HTML}{EAF4EF}
\definecolor{gold}{HTML}{DCA947}
\definecolor{sand}{HTML}{F7F3E8}
\definecolor{muted}{HTML}{5D6B66}
\definecolor{line}{HTML}{CFDDD6}

\hypersetup{colorlinks=true,urlcolor=green,linkcolor=green}
\pagestyle{fancy}
\fancyhf{}
\renewcommand{\headrulewidth}{0pt}
\fancyfoot[L]{\footnotesize\color{muted}Møtebrief · 27. august 2026}
\fancyfoot[R]{\footnotesize\color{muted}\thepage/2}
\setlength{\parindent}{0pt}
\setlength{\parskip}{3.2pt}
\setlist[itemize]{leftmargin=4.2mm,itemsep=1.4pt,topsep=2pt,parsep=0pt}
\titleformat{\section}{\large\bfseries\color{ink}}{}{0pt}{}[\vspace{-2pt}\color{gold}\rule{\linewidth}{0.7pt}]
\titleformat{\subsection}{\normalsize\bfseries\color{green}}{}{0pt}{}
\titlespacing*{\section}{0pt}{7pt}{4pt}
\titlespacing*{\subsection}{0pt}{5pt}{1pt}
\newcommand{\tagline}[1]{\colorbox{mint}{\parbox{0.965\linewidth}{\color{ink}\small #1}}}
\newcommand{\ask}[2]{\textbf{\color{green}#1}\quad #2}

\begin{document}

{\Huge\bfseries\color{ink} Møte med Anne Woie}\hfill
{\small\color{muted}\begin{tabular}{r}Næringssjef, Stavanger kommune\\Beslutnings- og samtalebrief\end{tabular}}

\vspace{2mm}
\tagline{\textbf{Mål for møtet:} Få en tydelig strategisk vurdering, riktig kommunal inngang og 2--4 konkrete introduksjoner som gjør en student--næringslivsturnering gjennomførbar som pilot i 2027.}

\section{Åpning: 30 sekunder}

``Vi utvikler en årlig, éndags møteplass der studenter og arbeidsgivere bygger relasjoner gjennom lagspill, middag og et strukturert sosialt program. Første pilot er planlagt for rundt 54 studenter og 18 bedriftsrepresentanter. Ambisjonen er å gjøre lokale karrieremuligheter mer synlige, skape varige relasjoner og gi flere studenter en grunn til å bli i eller vende tilbake til Stavanger-regionen. Vi ønsker din vurdering av strategisk treff, riktig forankring og hvem som bør sitte ved bordet videre.''

\section{Poengene som bør løftes}

\begin{tabularx}{\linewidth}{>{\raggedright\arraybackslash\bfseries\color{green}}p{34mm} X}
Samfunnsbehov & Regionen trenger å rekruttere og beholde kompetanse. Konseptet gjør arbeidsgivere og karriereveier konkrete tidlig, mens studentene fortsatt bygger tilhørighet.\\[2.2mm]
Annerledes møteplass & Flere timer på blandede lag gir mer naturlig og likeverdig kontakt enn korte standsamtaler. Middagsplassering og oppfølging gjør kontakten målrettet uten at dagen blir en karrieremesse.\\[2.2mm]
Bred regional verdi & Flere arbeidsgivere deltar på like vilkår. Tiltaket bør måles på relasjoner, praksis/intervjuer, lokal jobbinteresse og partnerfornyelse -- ikke bare antall deltakere.\\[2.2mm]
Realistisk pilot & Arbeidsscenarioet er 72 spillere og 18 lag. Budsjettet er om lag NOK 177,000; det første finansieringsmålet er NOK 140,000 i partnerskap slik at studentprisen kan holdes rundt NOK 800--900.\\[2.2mm]
Ansvarlig modell & Studentdata deles bare etter aktivt samtykke, til verifiserte virksomheter, for et tydelig formål og en avgrenset periode. Ingen bedrift får kontroll over deltakervalg eller garanterte rekrutteringsresultater.\\[2.2mm]
Repeterbar effekt & Start som én dokumentert pilot med før-/etterdata, partner- og studenttilbakemeldinger og tydelige kriterier for en årlig videreføring.\\
\end{tabularx}

\section{Hvorfor dette treffer hennes bord}

Stavangers planer vektlegger konkurransekraft, kompetanse, attraktiv studentby og samarbeid mellom kommune, kunnskapsmiljøer og næringsliv. Kunnskapsbyen peker uttrykkelig på å rekruttere og beholde arbeidskraft, bygge relasjoner med studenter og skape arbeidsrelevante koblinger. Anne kan derfor særlig vurdere strategisk relevans, regional forankring og kontaktveier -- men økonomisk støtte og formelle partnerskap må følge kommunens ordinære prosesser.

\vfill
{\small\color{muted}\textbf{Husk tonen:} Be om råd og presisjon, ikke en løs støtteerklæring. Gjør det enkelt å si ``ja, hvis \ldots'' og å navngi neste eier.}

\newpage
\small

{\LARGE\bfseries\color{ink} Hva Anne kan hjelpe med å avklare}

\section{Sju beslutninger å få ut av rommet}

\begin{enumerate}[leftmargin=7mm,itemsep=2.2pt,topsep=1pt]
\item \ask{Strategisk treff}{Hvilket kommunalt mål er den sterkeste inngangen: kompetanse, studentby, attraktivitet eller næringsutvikling? Hva må endres for at koblingen skal være troverdig?}
\item \ask{Geografi}{Kan Stavanger støtte eller inngå i en regional pilot dersom golfdelen skjer i Sola eller Sandnes, mens målgruppe, partnere og effekter er tydelig knyttet til Stavanger-regionen? Hvilke krav gjelder?}
\item \ask{Riktig eier}{Hvem i kommunen bør være fast kontakt videre, og hvilke saker hører hjemme hos næringsavdelingen, Kunnskapsbyen, Studentutvalget eller andre?}
\item \ask{Aktørkart}{Hvilke 2--4 introduksjoner gir størst fremdrift? Prioriter gjerne Kunnskapsbyen, UiS/BI/SiS, Næringsforeningen, studentorganisasjoner og relevante arbeidsgivernettverk.}
\item \ask{Kommunens mulige rolle}{Er den mest realistiske rollen faglig sparringspartner, døråpner, formidlingspartner, pilotpartner eller søkerveileder? Hva kan hun anbefale administrativt, og hva krever politisk eller formell behandling?}
\item \ask{Støtteløp}{Er næringsutviklingsmidler en reell mulighet for en pilot? Hvem er søknadsberettiget, hvilken 2026/2027-frist gjelder, og hvordan bør allmenn nytte, grønn profil, finansieringsmiks og måling dokumenteres?}
\item \ask{Legitimitet og avgrensning}{Hva må være på plass for at kommunen trygt kan omtale eller koble seg til tiltaket -- arrangørform, habilitet, åpenhet, personvern, alkoholhåndtering, partnerbalanse og bruk av kommunal merkevare?}
\end{enumerate}

\section{Vær tydelig på hva du ber om nå}

\rowcolors{2}{mint}{white}
\begin{tabularx}{\linewidth}{>{\bfseries}p{35mm} X p{38mm}}
\rowcolor{ink}\color{white}Område & \color{white}Ønsket utfall i møtet & \color{white}Bevis på fremdrift\\
Retning & En ærlig ``strategisk treff / treff med vilkår / feil spor'' & 2--3 konkrete justeringer\\
Forankring & Navngitt hovedkontakt og riktig kommunal kanal & Avtalt introduksjon eller oppfølging\\
Finansiering & Avklart om og hvordan støtteordningen er relevant & Kriterier, frist og søkerform\\
Nettverk & Prioritert liste over aktører og arbeidsgivere & 2--4 varme introduksjoner\\
Pilot & Enighet om hva som må dokumenteres før støtte/omtale & Kort kravliste og neste milepæl\\
\end{tabularx}

\section{Gode oppfølgingsspørsmål}

\begin{itemize}
\item ``Hva ville gjort at du selv fikk tillit til dette som et regionalt kompetansetiltak?''
\item ``Hvem vil savne en plass rundt bordet dersom vi går videre uten dem?''
\item ``Hvilken innvending vil komme først internt, og hvordan bør vi løse den nå?''
\item ``Hvis vi skal bevise effekt etter første år, hvilke tre måltall betyr mest for kommunen?''
\item ``Kan vi avslutte med navn, ansvar og én dato for neste steg?''
\end{itemize}

\section{Ikke lov for mye}

Ikke presenter kommunal finansiering, godkjenning eller promotering som sannsynlig før prosess og mandat er avklart. Ikke gjør tilskudd til hovedpoenget; be først om strategisk retning. Hold arbeidsgivertilgangen balansert, studentenes deling frivillig og den offentlige nytten tydelig. Dersom arenaen ligger utenfor Stavanger, løft geografien selv før Anne gjør det.

\section{Foreslått avslutning}

``Kan vi oppsummere dette i tre punkter: hvem vi følger opp, hva vi må sende dem, og hvilken beslutning vi skal være klare for neste gang?''

{\scriptsize\color{muted}
\textbf{Kildegrunnlag, kontrollert 27.08.2026:}
Stavanger kommune, \href{https://www.stavanger.kommune.no/naring-og-arbeidsliv/handlingsplan-for-naring-og-internasjonalt-samarbeid-2026-2027/}{Handlingsplan 2026--2027};
\href{https://www.stavanger.kommune.no/naring-og-arbeidsliv/kunnskapsbyenStavanger/temaplan-for-kunnskapsbyen/}{Temaplan for Kunnskapsbyen};
\href{https://www.stavanger.kommune.no/naring-og-arbeidsliv/midler-til-naringsutvikling/}{økonomisk støtte til næringsutvikling};
\href{https://www.stavanger.kommune.no/politikk/politiske-utvalg-i-stavanger/studentutvalget/}{Studentutvalget};
\href{https://www.stavanger.kommune.no/Stavangerbusinessregion/artikler/ons-2026/}{Stavanger Business Region / ONS 2026}.
Interne arbeidstall: prosjektbeskrivelse, partnerpakker og Bærheim-budsjett. Kommunens nettside viste ved kontrollen fortsatt eldre fristtekst for støtteordningen; aktuell frist må bekreftes i møtet.}

\end{document}
